\documentclass[a4paper]{article}
% Español (México) con XeLaTeX: fontspec para fuentes Unicode, polyglossia
% para la división silábica y los nombres del documento en la variante mexicana.
\usepackage{fontspec}
\usepackage{polyglossia}
\setdefaultlanguage[variant=mexican]{spanish}
% Los nombres chinos van como «pinyin (original)»: xeCJK da a los caracteres
% Han su propia fuente; sin ella XeLaTeX los omite con un aviso.
% FandolSong no cubre algunos caracteres de nombres (U+73FA); WenQuanYi Zen Hei sí.
\usepackage[AutoFallBack=true]{xeCJK}
\setCJKmainfont{FandolSong-Regular.otf}
\setCJKfallbackfamilyfont{\CJKrmdefault}{WenQuanYi Zen Hei}
% polyglossia no traduce los nombres de listings.
\AtBeginDocument{\providecommand{\lstlistingname}{}\renewcommand{\lstlistingname}{Listado}%
  \providecommand{\lstlistlistingname}{}\renewcommand{\lstlistlistingname}{Índice de listados}}
\usepackage{amsmath,amssymb}
\usepackage{graphicx}
\usepackage[margin=2.35cm]{geometry}
\usepackage[most]{tcolorbox}
\usepackage{etoolbox}
\usepackage{listings}
\usepackage{xcolor}
\usepackage{hyperref}
\usepackage{booktabs,longtable,tabularx,array}
\usepackage{subcaption}
\usepackage{float}
\usepackage{enumitem}
\usepackage{microtype}
\usepackage{fancyhdr}
\usepackage{needspace}

\definecolor{kbBack}{HTML}{F2F6FA}
\definecolor{kbFrame}{HTML}{9DB4CC}
\definecolor{kbTitle}{HTML}{52708F}
\definecolor{ibBack}{HTML}{FBF5E7}
\definecolor{ibFrame}{HTML}{C9A961}
\definecolor{ibTitle}{HTML}{7D5F1E}
\definecolor{wbBack}{HTML}{FCF2F0}
\definecolor{wbFrame}{HTML}{D6A096}
\definecolor{wbTitle}{HTML}{9E4F43}
\definecolor{dbBack}{HTML}{F4F7F3}
\definecolor{dbFrame}{HTML}{AEC2AC}
\definecolor{dbTitle}{HTML}{5F7F64}

\tcbset{softbox/.style={
  enhanced,breakable,boxrule=0.8pt,arc=2pt,left=3mm,right=3mm,
  top=2mm,bottom=2mm,coltitle=white,fonttitle=\bfseries,
  toptitle=0.6mm,bottomtitle=0.6mm
}}
\newtcolorbox{knowledgebox}[1]{softbox,colback=kbBack,colframe=kbFrame,colbacktitle=kbTitle,title={#1}}
\newtcolorbox{importantbox}[1]{softbox,colback=ibBack,colframe=ibFrame,colbacktitle=ibTitle,title={#1}}
\newtcolorbox{warningbox}[1]{softbox,colback=wbBack,colframe=wbFrame,colbacktitle=wbTitle,title={#1}}
\newtcolorbox{dialoguebox}[1]{softbox,colback=dbBack,colframe=dbFrame,colbacktitle=dbTitle,title={#1}}

\lstset{
  language=Python,basicstyle=\ttfamily\small,
  keywordstyle=\color{kbTitle}\bfseries,stringstyle=\color{wbTitle},
  commentstyle=\color{dbTitle},breaklines=true,frame=single,numbers=left,
  numberstyle=\tiny\color{gray},captionpos=b,columns=fullflexible,
  keepspaces=true,showstringspaces=false,extendedchars=true
}
\BeforeBeginEnvironment{lstlisting}{\Needspace{12\baselineskip}}

\hypersetup{colorlinks=true,linkcolor=kbTitle,urlcolor=kbTitle,citecolor=wbTitle}
\setlist{itemsep=0.25em,topsep=0.4em}
\setlength{\parindent}{2em}
\setlength{\parskip}{0.25em}
\newcolumntype{L}[1]{>{\raggedright\arraybackslash}p{#1}}

\providecommand{\notetitle}{Notas del curso: [tema del curso]}
\providecommand{\noteauthors}{Elaboradas a partir de los materiales públicos de Stanford CS25}
\providecommand{\notedate}{Versión reescrita en 2026}
\providecommand{\videochannel}{Stanford Online}
\providecommand{\videopublishdate}{2022}
\providecommand{\videoduration}{[duración del video]}
\providecommand{\videourl}{}
\providecommand{\videocoverpath}{cover.jpg}
\providecommand{\coursesite}{https://web.stanford.edu/class/cs25/}
\providecommand{\repourl}{https://github.com/hqhq1025/ai-course-notes}
\providecommand{\skillversion}{wdkns/wdkns-skills@39f1a04}

\newcommand{\figsource}[1]{\par\vspace{-0.3em}{\footnotesize\color{gray}\emph{Fuente: #1}}\par}
\newcommand{\teachervoice}[1]{\begin{knowledgebox}{Nota de clase}#1\end{knowledgebox}}
\newcommand{\term}[1]{\textbf{#1}}

\pagestyle{fancy}
\fancyhf{}
\fancyfoot[C]{\thepage}
\fancyfoot[R]{\footnotesize\color{gray}\href{\repourl}{AI Course Notes}}
\renewcommand{\headrulewidth}{0pt}
\renewcommand{\footrulewidth}{0pt}

\newcommand{\makecscover}{
\begin{titlepage}
\centering
\vspace*{0.7cm}
{\Large Stanford CS25 · Transformers United\par}
\vspace{0.8cm}
{\huge\bfseries \notetitle\par}
\vspace{0.6cm}
{\large \noteauthors\par}
\vspace{0.25cm}
{\large \notedate\par}
\vspace{0.9cm}
\ifdefempty{\videocoverpath}{}{
  \includegraphics[width=0.86\textwidth,height=0.43\textheight,keepaspectratio]{\videocoverpath}\par
}
\vfill
\begin{tcolorbox}[width=0.92\textwidth,colback=black!2!white,colframe=black!45,boxrule=0.7pt,arc=2pt]
\textbf{Autor o canal del video}: \videochannel\par
\textbf{Fecha de publicación}: \videopublishdate\par
\textbf{Duración del video}: \videoduration\par
\textbf{Sitio del curso}: \href{\coursesite}{\nolinkurl{\coursesite}}\par
\ifdefempty{\videourl}{}{\textbf{Enlace del video}: \href{\videourl}{\nolinkurl{\videourl}}\par}
\textbf{Normas de elaboración}: \skillversion\ + las normas source-first / teacher-voice / visual-QA de este repositorio
\end{tcolorbox}
\end{titlepage}
\tableofcontents
\newpage
}

\usepackage{multicol}

\renewcommand{\notetitle}{Lecture 32: Detrás del preentrenamiento de LLM: datos, gobernanza y evaluación de StarCoder}
\renewcommand{\noteauthors}{Elaborado a partir de la clase de Loubna Ben Allal, las 71 slides oficiales y los subtítulos revisados por personas}
\renewcommand{\notedate}{CS25 V4 · versión reescrita source-first de 2026}
\renewcommand{\videochannel}{Stanford Online}
\renewcommand{\videopublishdate}{Clase: 2024-05-23; publicación: 2024-06-07}
\renewcommand{\videoduration}{1:01:36}
\renewcommand{\videourl}{https://www.youtube.com/watch?v=jm2hyJLFfN8}
\renewcommand{\videocoverpath}{cover.jpg}

\newcommand{\lecturefigure}[3]{%
\begin{figure}[H]
  \centering
  \includegraphics[width=0.97\textwidth,height=0.47\textheight,keepaspectratio]{slides-images/#1}
  \caption{#2}
  \figsource{#3}
\end{figure}
}

\let\standardtableofcontents\tableofcontents
\renewcommand{\tableofcontents}{%
  \begingroup
  \small
  \setlength{\parskip}{0pt}
  \standardtableofcontents
  \endgroup
}

\begin{document}
\makecscover

\section{Fuente y pregunta central: detrás de un buen modelo no hay un solo comando de entrenamiento}

La versión anterior no incluía la canonical video URL, solo repetía la YouTube thumbnail como única figura de la clase y reducía el deck oficial de 71 páginas a una descripción breve. Esta clase se basa ahora en las Google Slides enlazadas directamente desde el sitio web personal de la ponente, en la grabación oficial de Stanford Online \texttt{jm2hyJLFfN8}, en los subtítulos oficiales \texttt{en-US} hechos por personas y en las primary sources de BigCode/StarCoder. Las 71 páginas de slides ya están guardadas en disco; 58 de ellas tienen contenido didáctico propio y las demás son solo separadores, estados intermedios de animaciones, marcas de finalización, páginas repetidas o agradecimientos.

\subsection{La clase no se reduce a la frase «los datos son lo más importante», sino que es una cadena de preguntas}

El título presenta StarCoder como use case, lo que significa que el objetivo de la clase no es repetir la configuración de un modelo, sino usar un modelo de código para responder preguntas generales de preentrenamiento: cuántos datos se necesitan, dónde encontrarlos, cómo filtrarlos, cómo gobernarlos, cómo evaluar y cuándo confiar en un benchmark. Cada capítulo posterior se desarrolla siguiendo esta estructura question--answer.

\lecturefigure{slide-01.jpg}{Behind the Scenes of LLM Pretraining: StarCoder use case}{Deck oficial de Loubna Ben Allal, página 1}

\lecturefigure{slide-02.jpg}{Ruta de la clase: responder «cómo entrenar un buen LLM» con una serie de preguntas}{Deck oficial de Loubna Ben Allal, página 2}

\begin{importantbox}{Las seis preguntas que esta clase debe responder}
\begin{enumerate}
  \item Con un compute dado, ¿cómo se reparten los parámetros del modelo y los tokens de entrenamiento?
  \item ¿Cómo se convierten las raw web/code/synthetic sources en un corpus entrenable?
  \item ¿Cómo se valida la intuición detrás de un filter mediante ablation con modelos pequeños?
  \item ¿Cómo entran la license, el opt-out, la PII y la decontamination al pipeline?
  \item ¿Cómo se da formato de sequence de entrenamiento a un repository/file/commit/notebook?
  \item ¿Cómo evita un leaderboard la contamination, la dependencia de un solo benchmark y la fuga temporal?
\end{enumerate}
\end{importantbox}

\teachervoice{00:00:50--00:02:10, Ben Allal dice que toda la exposición tiene una sola pregunta, «muy loaded», y que las slides irán planteando subpreguntas para luego responderlas. StarCoder es el caso que lleva los principios abstractos a un proyecto real.}

\subsection{Avances de los Open-model: qué capacidades prácticas trae abrir los pesos}

El memo ``No Moat'', que circuló en 2023, representa un juicio: que los modelos cerrados vayan adelante no significa que los abiertos nunca puedan alcanzarlos. La clase no lo cita para declarar terminada la competencia, sino para explicar por qué publicar los weights genera un nuevo ecosistema de deployment/fine-tuning.

\lecturefigure{slide-03.jpg}{El juicio «No moat»: la comunidad abierta completa poco a poco las piezas del entrenamiento}{Deck oficial de Loubna Ben Allal, página 3}

Una vez abiertos los weights, la quantization puede reducir el memory footprint, un runtime local permite desplegar en consumer hardware y el parameter-efficient fine-tuning permite personalizar con los datos de una organización. Este valor no equivale a que el modelo alcance la capacidad de los sistemas cerrados en todas las tareas, sino a que \textbf{el espacio de opciones verificables, modificables y desplegables de forma privada es mayor}.

\lecturefigure{slide-04.jpg}{Ejemplo de la clase de 2024: un modelo con evaluaciones cercanas a GPT-4 puede ejecutarse en un equipo de escritorio de consumo de gama alta}{Deck oficial de Loubna Ben Allal, página 4}

\begin{warningbox}{Un «rendimiento cercano» debe ir ligado a una versión, una métrica y una fecha}
La slide es una instantánea de mayo de 2024. Arena, MMLU, coding, seguridad, contexto largo, tool use, latency y cost dan clasificaciones distintas; esta clase no convierte la cercanía en una tabla concreta en una parity general.
\end{warningbox}

\teachervoice{00:02:10--00:03:20, la ponente enfatiza que el beneficio directo de los open weights es la quantization, el consumer deployment y el fine-tuning comunitario, no una frase abstracta como «lo abierto es mejor».}

\subsection{Las open releases aumentan, pero la transparencia del entrenamiento sigue siendo insuficiente}

Modelos públicos como Gemma y Mistral amplían el open ecosystem; la gráfica de Arena también muestra que la open/closed gap se redujo entre 2023 y 2024. Al leer estas dos páginas, primero hay que ver qué pregunta responde cada una: la primera trata del crecimiento de la oferta; la segunda, de la tendencia bajo un protocolo específico de evaluación de preferencias humanas.

\lecturefigure{slide-05.jpg}{Aumentan los lanzamientos de modelos abiertos: Gemma, Mistral y otros entran a la competencia}{Deck oficial de Loubna Ben Allal, página 5}

\lecturefigure{slide-06.jpg}{Perspectiva de Arena en May 2024: se reduce la brecha entre modelos open/closed}{Deck oficial de Loubna Ben Allal, página 6}

\begin{knowledgebox}{Qué puede demostrar la evidencia de Arena}
Arena forma un Elo-like ranking a partir de la pairwise preference de los usuarios, lo que sirve para medir la preferencia por respuestas interactivas; por sí solo no puede medir la legalidad de los datos de entrenamiento, la confiabilidad factual, la rare-language coverage, la code execution correctness ni el costo de despliegue. La tendencia sirve como referencia, pero los límites de capacidad siguen requiriendo una evaluación multidimensional.
\end{knowledgebox}

\subsection{Por qué se publican los pesos pero no los datos ni la recipe}

La ponente da dos razones prácticas: primero, la disclosure puede provocar license/copyright scrutiny; segundo, los equipos quieren conservar su competitive edge. Así, los weights del modelo pueden descargarse, pero no es posible saber de dónde vienen las muestras, cómo se eligieron los filters, cómo se asignó el número de tokens ni si hubo fuga en la evaluation.

\lecturefigure{slide-07.jpg}{Limitaciones de los modelos públicos: los datos y los detalles de entrenamiento suelen faltar}{Deck oficial de Loubna Ben Allal, página 7}

\begin{warningbox}{La falta de reproducibilidad no es «un poco menos de documentación»}
Si faltan la dataset version, el dedup threshold, los sampling weights, el tokenizer, los training tokens, el optimizer schedule, la benchmark decontamination y las seeds, los investigadores no pueden determinar si el rendimiento proviene de la architecture, de los data o de una fuga en la evaluación, ni pueden reutilizar de forma confiable las decisiones de gobernanza.
\end{warningbox}

\teachervoice{00:04:10--00:05:05, la ponente expone abiertamente las razones legales y competitivas para no divulgar. No por ello renuncia a resumir regularidades públicas, sino que enfatiza la necesidad de reunir evidence a partir de varias releases públicas.}

\subsection{Model, GPUs, Data: de los tres vértices del triángulo, esta clase elige data}

Transformer, Mamba y MoE pueden ser opciones de model, y las GPU fijan el límite del budget; pero con architecture y training technique similares, la calidad y la mixture de los datos suelen ser la principal diferencia dentro de un presupuesto fijo. Aquí, el «data backbone» es el focus de la clase; no se afirma que la architecture y el hardware sean irrelevantes.

\lecturefigure{slide-11.jpg}{Los tres tipos de recursos para entrenar un buen LLM: Model, GPUs y Data}{Deck oficial de Loubna Ben Allal, página 11}

\begin{knowledgebox}{Escribir los tres como una optimización con restricciones}
El objetivo de entrenamiento puede escribirse de forma aproximada como
\[
\min_{\theta,\mathcal{D},r}\ L(\theta;\mathcal{D},r)
\quad\text{s.t.}\quad
C_{\text{train}}\le B,\ C_{\text{serve}}\le S,
\]
donde $\mathcal{D}$ son los datos y la mixture, $r$ es la recipe de entrenamiento y $B,S$ son, respectivamente, los presupuestos de entrenamiento y de servicio. La clase optimiza principalmente $\mathcal{D}$, pero las restricciones siguen viniendo del modelo y de las GPU.
\end{knowledgebox}

\teachervoice{00:05:05--00:06:20, la ponente dice que las architecture/techniques se parecen cada vez más, por lo que, «con un presupuesto dado», son los datos los que marcan la diferencia entre modelos. La condición es given budget, no que a cualquier escala solo importen los datos.}

\subsection{Resumen de la sección}

Los open weights amplían el espacio de despliegue y personalización, pero la transparencia del entrenamiento sigue siendo la principal carencia. Por eso la clase reduce el preentrenamiento a una labor de ingeniería sujeta a la vez a restricciones de compute, serving, data, governance y evaluation, y StarCoder ofrece un caso de extremo a extremo que se puede verificar.
\section{Cuántos datos se necesitan: Scaling law y costo del ciclo de vida}

Antes de decidir las fuentes de datos, primero hay que estimar cuántos tokens se necesitan. El valor de la Scaling law es ofrecer curvas empíricas para repartir el presupuesto; su limitación es que el optimum de la loss de entrenamiento no coincide automáticamente con el optimum del costo de despliegue, de las restricciones de datos ni del domain-specific performance.

\subsection{Tres preguntas: cuántos, de dónde y cómo filtrar}

La sección anterior planteó el modelo, las GPU y los datos como un mismo problema con restricciones; esta sección empieza a desglosar la variable de los datos. Al leer la figura, las tres preguntas deben verse como un pipeline con dependencias: la cantidad de datos determina la escala de la collection, la fuente determina la provenance y el domain, y el filter determina la quality/diversity; las tres interactúan entre sí, y no se puede fijar una primero y tratar las demás como simples pasos de limpieza.

\lecturefigure{slide-13.jpg}{Las tres preguntas sobre los datos de entrenamiento: how much, where, how to filter}{deck oficial de Loubna Ben Allal, página 13}

\begin{importantbox}{Hay que distinguir raw data, retained data y training tokens}
Los raw bytes son el volumen recolectado; los retained bytes son los archivos que quedan después de filter/dedup; los unique tokens son el conteo sin duplicados después del tokenizer; los training tokens incluyen además los epochs/repeats y el mixture sampling. Si un reporte dice «3T tokens» sin aclarar a qué nivel se refiere, no se puede reproducir.
\end{importantbox}

\subsection{Scaling law: repartir el compute entre parámetros y tokens}

Determinar «cuántos datos se necesitan» no puede basarse solo en la capacidad de disco disponible; hay que colocar al mismo tiempo el tamaño del modelo, los tokens de entrenamiento y el presupuesto dentro de la curva empírica. Esta sección presenta primero la forma de la scaling-law y luego usa las conclusiones distintas de Kaplan y de Chinchilla para mostrar que el llamado optimum es una asignación de recursos obtenida bajo un conjunto de supuestos y un fitted range, no una constante independiente de los datos y del hardware.

La empirical law de estilo Kaplan suele escribirse como
\[
L(N,D)=L_\infty+aN^{-\alpha}+bD^{-\beta},
\]
donde $N$ es el número de parámetros, $D$ son los tokens de entrenamiento y $L_\infty$ es la loss irreducible. Para un Transformer dense, los FLOPs de entrenamiento suelen aproximarse como $C\approx 6ND$; con $C$ fijo, hay que elegir la combinación de modelo y datos.

\lecturefigure{slide-14.jpg}{How much data: usar scaling laws para repartir el compute budget}{deck oficial de Loubna Ben Allal, página 14}

\lecturefigure{slide-15.jpg}{Kaplan y Chinchilla: conclusiones compute-optimal distintas y ejemplos}{deck oficial de Loubna Ben Allal, página 15}

Los primeros resultados de Kaplan favorecían modelos más grandes con menos tokens; Chinchilla, con experimentos más amplios, los corrigió hacia un crecimiento más equilibrado entre parámetros y tokens. La slide compara GPT-3 175B con alrededor de 300B tokens y Chinchilla 65B con alrededor de 1.4--1.6T tokens; la conclusión didáctica es: \textbf{con el mismo compute, los modelos anteriores podían tener demasiados parámetros y estar subentrenados}.

\begin{warningbox}{Una razón constante no es una ley permanente}
«Alrededor de 20 tokens por parámetro» depende de la familia de modelos, los datos, el optimizer, el context y el fitted range. Los modelos posteriores suelen hacer overtrain de forma intencional (intentionally overtrain), y los cambios de domain o de data quality también desplazan el optimum.
\end{warningbox}

\teachervoice{00:06:20--00:09:10, la ponente usa Kaplan/Chinchilla para explicar que la scaling law es una herramienta de asignación de presupuesto, y advierte que estas conclusiones provienen de artículos y datos específicos, por lo que no conviene memorizar un único número de token-per-parameter.}

\subsection{Cómo leer las curvas de Chinchilla: cada compute budget tiene un mínimo}

La sección anterior presentó la fórmula abstracta; esta pasa a las curvas empíricas para mostrar cómo aparece el optimum en la gráfica. Al leer la figura, fije primero un color: el eje horizontal es el número de parámetros y el vertical es la validation loss; cada curva primero baja y luego sube, porque un modelo demasiado pequeño carece de capacidad y uno demasiado grande, con presupuesto fijo, no recibe suficientes tokens; el mínimo es el training-compute optimum para ese budget.

\lecturefigure{slide-16.jpg}{Chinchilla empirical curves: el mínimo en número de parámetros con compute fijo}{deck oficial de Loubna Ben Allal, página 16}

Si $C=6ND$, al sustituir $D=C/(6N)$ en la law empírica se pueden obtener las leyes de potencias de $N^*(C)$ y $D^*(C)$. Un entrenamiento real no se limita a resolver la fórmula: también debe considerar el discrete hardware shape, la parallel efficiency, el batch, el context y la disponibilidad de datos.

\begin{knowledgebox}{Cuatro cosas que la curva no dice}
Por lo general no contabiliza el serving cost; no garantiza que las downstream tasks mejoren al mismo ritmo; no describe las license/quality constraints; y tampoco incluye el costo de las fallas de entrenamiento (failure), la checkpoint selection ni la hyperparameter search.
\end{knowledgebox}

\subsection{Tres líneas de scaling evidence: Compute, Data y Model}

Las curvas de Chinchilla muestran el mínimo dentro de un presupuesto fijo; esta página va más allá y grafica la loss por separado contra el compute, el data size y el model size, mostrando los intervalos donde se ajusta una recta. Al leer la figura, confirme primero si las demás variables están cerca del óptimo y revise el rango del eje horizontal; de lo contrario, aumentar un solo eje puede limitarse a empujar el sistema a un régimen de under-training o data-limited.

\lecturefigure{slide-17.jpg}{Scaling evidence: variación de la loss como ley de potencias del compute, los data y el model size}{deck oficial de Loubna Ben Allal, página 17}

\begin{warningbox}{Las tres rectas no se pueden combinar en una causalidad incondicional}
Estas gráficas son relaciones marginales empíricas. Si la data quality baja con la escala, si la tasa de repetición aumenta o si el modelo sale del intervalo de estabilidad de la optimización, la misma pendiente no se prolonga indefinidamente. Toda predicción debe reportar el fitted range y la confidence.
\end{warningbox}

\subsection{De Chinchilla al overtraining: qué cambió en el optimum}

La línea de tiempo muestra la tendencia de Chinchilla 65B/1.4T, LLaMA 7B/1T, LLaMA-3 8B/15T, entre otros. Los equipos entrenan deliberadamente modelos más pequeños con más tokens, porque el costo adicional de un solo entrenamiento se compensa con menor inference memory, latency y deployment complexity a largo plazo.

\lecturefigure{slide-18.jpg}{Scaling-law timeline: de compute-optimal al overtraining de modelos pequeños}{deck oficial de Loubna Ben Allal, página 18}

Se define el costo del ciclo de vida del modelo como
\[
C_{\text{life}}=C_{\text{train}}+Q\cdot C_{\text{infer/token}}+C_{\text{memory}}+C_{\text{ops}},
\]
donde $Q$ es el total de tokens durante el periodo de servicio. Si $Q$ es muy grande, entrenar más un smaller model puede reducir el costo total.

\teachervoice{00:09:10--00:11:10, la ponente dice que la Chinchilla scaling law no incluye el inference cost en su objetivo; por eso hoy se eligen modelos más pequeños entrenados por más tiempo, pagando una vez un costo mayor durante el entrenamiento.}

\subsection{De GPT-1 a GPT-4: la scale amplía tanto la factura del entrenamiento como la del servicio}

La discusión anterior sobre el optimum sigue siendo abstracta; aquí se usan los cambios de orden de magnitud de la serie GPT para poner en un mismo cuadro la factura del entrenamiento y la del servicio. Las slides muestran el crecimiento de escala mediante dataset/model size y costos de entrenamiento aproximados; las cifras son estimaciones públicas, no registros contables auditados. Su función didáctica es recordar al lector que aumentar los parámetros no solo afecta a un entrenamiento, sino que encarece a largo plazo la latency por token, la GPU memory y el servicio concurrente.

\lecturefigure{slide-19.jpg}{De GPT-1 a GPT-4: crecimiento del dataset, el model size y el costo de entrenamiento}{deck oficial de Loubna Ben Allal, página 19}

\lecturefigure{slide-20.jpg}{Inference gets more expensive: número de parámetros, memoria de GPU y necesidad de GPU}{deck oficial de Loubna Ben Allal, página 20}

Si la precisión de los pesos es de $b$ bytes, la memory de una sola copia de los pesos es al menos
\[
M_{\text{weights}}=bN,
\]
y además se requieren el KV cache, las activations, el workspace y la redundancia del paralelismo. La cuantización reduce $b$, pero no reduce automáticamente todos los cuellos de botella (bottleneck) de compute y bandwidth.

\begin{knowledgebox}{Por qué se prefieren code models pequeños en el IDE}
El autocompletado de código exige baja latency, alta concurrencia y un repository context relativamente largo; cada pulsación de tecla del usuario puede generar una solicitud. Un modelo con mejor benchmark pero mayor latencia no es necesariamente el optimum del producto; esto también explica por qué StarCoder2 se ofrece en varios tamaños.
\end{knowledgebox}

\subsection{Compute-optimal no equivale a optimal: la advertencia de la Harm's law}

La sección anterior mostró que los modelos grandes pagan el inference cost una y otra vez; por eso esta sección cambia el objetivo de la loss de un solo entrenamiento al costo del ciclo de vida. Las dos slides señalan que entrenar por más tiempo un smaller model eleva el upfront compute, pero puede reducir el costo de cada deployment; la llamada ``Harm's law'' es solo un nombre humorístico para el costo adicional de entrenar más allá del punto de Chinchilla, no una nueva ley natural.

\lecturefigure{slide-21.jpg}{Compute-optimal $\neq$ deployment-optimal: por qué sobreentrenar modelos pequeños}{deck oficial de Loubna Ben Allal, página 21}

\lecturefigure{slide-22.jpg}{Harm's law: el compute adicional de entrenar más allá del punto de Chinchilla}{deck oficial de Loubna Ben Allal, página 22}

\begin{importantbox}{Para decidir, trace al menos dos curvas}
Una es training loss versus training compute; la otra es end-to-end product cost/quality versus model size. Si solo se optimiza la primera, se pasan por alto la repeated inference, la memory capacity, la energy, el edge deployment y la user latency.
\end{importantbox}

\subsection{Data constraint y quality desplazan la scaling curve}

La página de lecturas adicionales introduce el data-constrained scaling y DeepSeek LLM: cuando los unique data son limitados, repetir el entrenamiento cambia los beneficios; un corpus de mayor calidad también cambia la model/data allocation. Dicho de otro modo, $D$ no es solo un conteo de tokens: también conlleva una distribution y una quality.

\lecturefigure{slide-23.jpg}{Lecturas adicionales sobre scaling laws: data constraint y quality-dependent allocation}{deck oficial de Loubna Ben Allal, página 23}

\begin{warningbox}{El domain scaling sigue siendo un problema underexplored}
En la sesión de Q\&A, la ponente dijo que Chinchilla tiene cierto respaldo entre English/code, pero que en domains como el medical puede ser distinto; incluso dentro del mismo domain, cambiar a datos más curated puede desplazar el optimum. No copie directamente un generic web fit a un campo especializado.
\end{warningbox}

\teachervoice{00:11:10--00:13:20, la ponente presenta la data quality y el domain como caveats de la scaling-law, y señala que los resultados de DeepSeek muestran que cambiar el conjunto de datos dentro del mismo campo también modifica la asignación.}

\subsection{Resumen de la sección}

La Scaling law responde a «cómo repartir con un compute de entrenamiento fijo», pero no abarca todo el ciclo de vida del modelo. El optimum real depende también del serving volume, la memory, la latency, la unique-data constraint, el domain y la quality; por eso los equipos actuales suelen elegir modelos más pequeños con más tokens.
\section{De dónde vienen los datos: Web, Code, Curated y Synthetic}

Una vez fijado el token budget, la segunda pregunta es la source. Las distintas fuentes no solo difieren en contenido: también implican distintos costos de rastreo, licencia, deduplicación, formato y verification. La clase divide las fuentes en web data, GitHub/code, curated sources y synthetic data; StarCoder, por su parte, muestra cómo convertir un code corpus en un producto de datos que puede inspeccionarse.

\subsection{Source taxonomy: la fuente determina los costos de gobernanza posteriores}

La sección anterior solo estimó el token budget; esta sección empieza a responder de dónde vienen esos tokens, porque la source determina los costos posteriores de filter, license y evaluation. La web ofrece amplia cobertura, pero con mucho ruido; GitHub tiene estructura y metadatos de licencia, pero contiene duplicados y secrets; las curated sources son de alta calidad, pero de escala limitada; la synthetic data puede generarse de forma dirigida, pero depende del teacher, el prompt y la seed diversity.

\lecturefigure{slide-26.jpg}{Las cuatro categorías de fuentes de LLM training data}{Deck oficial de Loubna Ben Allal, página 26}

\begin{knowledgebox}{Las fuentes no son etiquetas mutuamente excluyentes}
Los repositorios de código contienen README/Markdown, las páginas web en lenguaje natural contienen código, y los synthetic prompts a menudo se derivan de web/curated documents. La mixture real debe registrarse según la provenance y la transformation lineage, no solo con los porcentajes finales por domain.
\end{knowledgebox}

\subsection{Common Crawl: una escala tan grande que obliga a filtrar antes de descargar/procesar}

Common Crawl ofrece instantáneas periódicas de la web; el raw corpus puede alcanzar cientos de TB. Tokenizarlo directamente no solo es costoso, sino que conserva boilerplate, spam, machine-generated pages, duplicates, PII y texto de baja calidad; por eso, normalmente se reutiliza un filtered web dataset existente o se aplica distributed filtering.

\lecturefigure{slide-27.jpg}{Web data: la escala de Common Crawl y la carga de filtrado}{Deck oficial de Loubna Ben Allal, página 27}

\begin{warningbox}{Rastrear una página web no equivale a tener permiso para entrenar con ella}
La accesibilidad, robots, copyright, terms, personal data y jurisdiction son cuestiones distintas. Un pipeline técnico puede registrar la procedencia y retirar muestras, pero no puede resolver automáticamente todos los juicios legales y éticos.
\end{warningbox}

\teachervoice{00:13:20--00:16:20, la ponente subraya la escala original de Common Crawl y el filtering cost, y recomienda entender el procesamiento de los datasets existentes en lugar de tratar el raw crawl como datos listos para usarse.}

\subsection{FineWeb: un mejor filter se demuestra con la early-training curve}

La slide de FineWeb compara las curvas de entrenamiento de modelos pequeños sobre varios web corpora. Al leer la figura, primero hay que confirmar que la architecture, el token budget, la evaluation y la seed sean los mismos; después, observar qué corpus alcanza antes un score más alto. Lo que la curva respalda es que el filter pipeline supera al control en ese entorno, no que sea el mejor para cualquier modelo/tarea.

\lecturefigure{slide-29.jpg}{FineWeb ablation: early-training performance de distintos web corpora}{Deck oficial de Loubna Ben Allal, página 29}

\begin{importantbox}{La operational definition de data quality}
Alta calidad no significa «parecer un buen artículo». En el preentrenamiento, debe manifestarse como menor validation loss, downstream learning más rápido, menos memorization/toxicity/PII y mejor domain coverage, y debe sostenerse con matched compute y varias seeds.
\end{importantbox}

\subsection{The Stack: los datos de código requieren raw collection, inspection y opt-out}

El código no es como una página web de texto plano: los archivos tienen estructura (repository, path, license, language, stars, commit history, etc.) y también pueden incluir generated files, vendor copies, secrets y benchmark solutions. Por eso, el pipeline de The Stack no solo recolecta: también realiza detección de lenguaje, gestión de license/metadata, dedup e inspection por parte de la comunidad.

\lecturefigure{slide-30.jpg}{Code data: panorama de collection y processing de The Stack}{Deck oficial de Loubna Ben Allal, página 30}

\lecturefigure{slide-31.jpg}{The Stack v1: GitHub collection, raw data, dedup, inspection y opt-out}{Deck oficial de Loubna Ben Allal, página 31}

\begin{knowledgebox}{El lugar del opt-out en el pipeline}
El opt-out requiere stable repository/file identifiers, herramientas públicas de consulta, una lista versionada de las solicitudes procesadas y un mecanismo de exclusión para el future retraining. No consiste en borrar un enlace web una vez terminado el entrenamiento, sino que forma parte de la data lineage.
\end{knowledgebox}

\teachervoice{00:17:10--00:19:10, la ponente explica que BigCode publicó con la v1 una inspection tool y un opt-out form, para que los desarrolladores pudieran revisar el código y solicitar su exclusión de entrenamientos futuros; este es un diseño central de los proyectos de datos abiertos.}

\subsection{The Stack v2: Software Heritage y más datos del proceso de desarrollo}

La v2 ya no depende principalmente de clonar GitHub de forma directa: extrae los datos del Software Heritage archive e incorpora issues, pull requests, commits, Jupyter/Kaggle notebooks y math/coding datasets. Así, el corpus no solo enseña el código final, sino que también contiene los problemas, las modificaciones, las discusiones y los entornos de ejecución.

\lecturefigure{slide-32.jpg}{The Stack v2: Software Heritage, filtrado y datos de desarrollo de múltiples fuentes}{Deck oficial de Loubna Ben Allal, página 32}

\lecturefigure{slide-33.jpg}{The Stack v1 y v2: raw size, filtered size, languages y tokens}{Deck oficial de Loubna Ben Allal, página 33}

\begin{importantbox}{La brecha entre raw scale y usable scale es trabajo de ingeniería}
El raw size de la v2 es aproximadamente diez veces el de la v1, y después del filtrado sigue siendo unas cuatro a cinco veces el de la v1. La gran mayoría de lo eliminado no es «desperdicio», sino duplicates, contenido de baja calidad, unsupported formats, PII o lenguajes no adecuados; la tasa de filtrado por sí sola no indica si el resultado es bueno o malo: hay que observar la diversity conservada y los resultados del modelo.
\end{importantbox}

\teachervoice{00:19:10--00:20:20, la ponente sitúa las mejoras centrales de la v2 en la archive source, más lenguajes y datos auxiliares de desarrollo, y no solo en «más tokens».}

\subsection{Synthetic data: de Phi a un corpus de generación controlable}

La serie Textbooks Are All You Need/Phi muestra que, con datos de alta calidad con estilo de libro de texto generados por GPT-3.5/4, es posible entrenar modelos pequeños potentes con un corpus mucho menor que uno de web-scale. Lo que esto suscitó no fue la conclusión de que «lo synthetic sustituye a los datos reales», sino que generator, prompt, seed y verification se convirtieran en nuevas variables de la ingeniería de datos.

\lecturefigure{slide-34.jpg}{Caso catalizador de la synthetic data: Textbooks Are All You Need}{Deck oficial de Loubna Ben Allal, página 34}

Sea $s$ la seed/context, $p$ el prompt template y $G$ el generator; la synthetic sample es
\[
y\sim G(\cdot\mid p,s).
\]
La calidad de los datos depende de la cobertura de $s$, las restricciones de $p$, la capacidad de $G$ y el verifier a posteriori $v(y)$; no puede depender solo de la cantidad de muestras.

\begin{warningbox}{Synthetic feedback loop}
Si el generator y el verifier comparten puntos ciegos, la generación repetida amplifica la homogeneización del estilo, los errores factuales y el representation collapse. Conviene registrar teacher/version, seed provenance, rejection reasons, novelty y human spot checks.
\end{warningbox}

\subsection{Cosmopedia: ampliar la diversidad con web/curated seeds}

Cosmopedia usa Mixtral-8x7B para generar unos 25B tokens de libros de texto, blogs, cuentos, etc. La técnica central no es una simple instrucción del tipo «haz que el modelo escriba un libro de texto», sino extraer context de web samples, Stanford courses, WikiHow, etc., y colocar en el prompt el topic junto con una seed concreta, para limitar la hallucination y aumentar la dispersión temática.

\lecturefigure{slide-35.jpg}{Cosmopedia: synthetic text dataset a gran escala generado con modelos abiertos}{Deck oficial de Loubna Ben Allal, página 35}

\lecturefigure{slide-36.jpg}{Composición de los datos de Cosmopedia: web seeds, curated sources y múltiples géneros}{Deck oficial de Loubna Ben Allal, página 36}

\begin{knowledgebox}{Las cuatro auditorías de un synthetic corpus}
\begin{enumerate}
  \item \textbf{Coverage:}¿están dispersos el topic, el language, la difficulty y el style?
  \item \textbf{Grounding:}¿la generación está restringida por la seed/context?
  \item \textbf{Verification:}¿se revisan de forma automática/humana los hechos, el código y las matemáticas?
  \item \textbf{Lineage:}¿pueden rastrearse teacher, prompt, seed, filter y license?
\end{enumerate}
\end{knowledgebox}

\teachervoice{00:20:20--00:24:20, la ponente dice que la synthetic data es muy reciente y muy prometedora, pero que obtener diversity es «very tricky». Cosmopedia usa una gran cantidad de web/curated seeds, en lugar de ampliar sin límite a partir de unos pocos prompts vagos.}

\subsection{Resumen de la sección}

La source de los datos determina los costos de procesamiento posteriores. La web aporta cobertura; el code, estructura; las curated sources, una signal alta; y la synthetic data, ampliación dirigida. Cualquier fuente debe ir acompañada de un diseño de provenance, filter, license/opt-out y evaluation.
\section{Cómo encontrar un buen Filter: no confíes en la intuición, haz primero una Ablation}

«Limpiar datos» es lo que con más facilidad se reduce a una lista de reglas, pero lo verdaderamente difícil es juzgar si una regla mejora el modelo. El curso convierte el filter discovery en ciencia experimental: la manual inspection propone candidatos, la small-model ablation estima el efecto causal, y el uso de varias seeds y de high-signal benchmarks reduce el ruido.

\subsection{De la afirmación sobre la calidad a un pipeline ejecutable}

La primera evidence slide cita la postura de Yi/Open Foundation Models: una standard architecture con datos de alta calidad también puede rendir muy bien. Esa slide solo aporta la motivación; la página siguiente muestra stages ejecutables como language filtering, quality rules, dedup y semantic/topic.

\lecturefigure{slide-39.jpg}{Por qué filtrar: el efecto de la data de alta calidad sobre una architecture estándar}{Deck oficial de Loubna Ben Allal, página 39}

\lecturefigure{slide-40.jpg}{Pipeline general de data-cleaning para pretraining}{Deck oficial de Loubna Ben Allal, página 40}

Un rule/perplexity filter común puede escribirse como una condición de conservación
\[
\mathbb{1}\{q(x)\in[a,b]\}\cdot \mathbb{1}\{\operatorname{PPL}_{r}(x)<\tau\},
\]
donde $q(x)$ es una heuristic como la longitud, las líneas repetidas o la proporción de símbolos, y $r$ es el reference model. Los umbrales $a,b,\tau$ deben calibrarse con muestras y con ablation.

\begin{warningbox}{El doble filo del perplexity filter}
Una PPL alta puede indicar basura, pero también un idioma poco común, código, fórmulas o conocimiento nuevo; una PPL baja puede indicar plantillas y repetición. Usar la PPL de un reference model introduce su bias en el dataset, por lo que debe combinarse con métricas de diversity/coverage.
\end{warningbox}

\subsection{En busca de la mejor técnica de filtrado: la manual inspection es solo el punto de partida}

La sección anterior enumeró los stages de language, rule, perplexity y dedup, pero aún no respondió cómo elegir los umbrales. Esta sección convierte el diseño de filters en un experimento refutable: primero se examinan los datos y se explican las anomalías, y después se entrenan modelos pequeños con matched compute para verificarlo; para la evaluación se eligen benchmarks que ya dan signal en el early training, y se usan distintas seeds para no confundir el optimizer noise con una filter gain.

\lecturefigure{slide-42.jpg}{Filter discovery: inspección manual, ablation con modelos pequeños, benchmarks de alta signal y múltiples seeds}{Deck oficial de Loubna Ben Allal, página 42}

Si la diferencia de puntuación entre el filter $f$ y el baseline es $\Delta_s$, la media y el error estándar a lo largo de $k$ seeds pueden escribirse como
\[
\bar\Delta=\frac{1}{k}\sum_{s=1}^{k}\Delta_s,
\qquad
\operatorname{SE}(\bar\Delta)=\frac{\operatorname{sd}(\Delta_s)}{\sqrt{k}}.
\]
Solo cuando el effect supera el noise vale la pena aplicar la regla a un TB-scale corpus.

\begin{lstlisting}[caption={Protocolo experimental mínimo de una ablation de filtrado de datos}]
for filter_candidate in candidates:
    subset = apply_filter(raw_subset, filter_candidate)
    for seed in seeds:
        model = train_small_model(subset, matched_tokens, seed)
        scores[filter_candidate, seed] = evaluate(high_signal_suite)
compare_mean_variance_cost_and_retained_diversity(scores)
\end{lstlisting}

\teachervoice{00:27:10--00:30:00, la ponente presenta dos resultados negativos: filtrar por comment density casi no produjo mejora; eliminar los repositorios con menos de 5 stars retiró más del 70\% de los datos y dio la peor ablation. La intuición debe ceder ante el experimento.}

\subsection{Las 200+ ablations de FineWeb: la recipe de datos también necesita un search budget}

El equipo de FineWeb entrenó más de 200 ablation models de alrededor de 1B de parámetros y alrededor de 30B tokens para elegir los filters. Esta escala muestra que la data research no es un preprocesamiento barato; requiere model-training budget, un benchmark harness y disciplina estadística.

\lecturefigure{slide-43.jpg}{FineWeb: búsqueda de la filtering recipe con 200+ ablation models}{Deck oficial de Loubna Ben Allal, página 43}

\begin{importantbox}{Tratar la data recipe como una hyperparameter search}
Filter type, threshold, dedup radius, language mixture, source weights y sampling temperature son todos hiperparámetros. Igual que en la architecture search, hay que evitar el test-set tuning, conservar una held-out evaluation y registrar los experimentos fallidos.
\end{importantbox}

\subsection{Resumen de la sección}

Un buen filter no es aquel cuya regla parece razonable, sino aquel que mejora el modelo con matched compute, múltiples seeds y varios early-signal benchmarks, y que al mismo tiempo preserva los objetivos de diversity y governance. El fracaso del star filtering es una de las lecciones prácticas más importantes de esta clase.
\section{De The Stack a StarCoder: pipeline de extremo a extremo para datos de código}

Esta sección lleva el filtering protocol general a un code corpus real. La particularidad de los datos de código proviene de las diferencias entre lenguajes, la duplicación de repositorios, los generated files, los secrets, las benchmark solutions, los metadata y los eventos de desarrollo estructurados. Cada paso cambia las tareas que el modelo puede aprender.

\subsection{De raw a training data: una gran reducción no equivale a desperdicio}

The Stack v1 pasa de unos 6.4TB de source code a los cerca de 800GB que usa StarCoder; v2 pasa de unos 32.1TB a unos 3.6TB. Esta reduction acumula la selección de lenguajes, las quality rules, el dedup, la PII y las restricciones de formato.

\lecturefigure{slide-44.jpg}{StarCoder/The Stack v1: de 6.4TB raw a 800GB de training data}{Deck oficial de Loubna Ben Allal, página 44}

\lecturefigure{slide-45.jpg}{StarCoder2/The Stack v2: de 32.1TB raw a 3.6TB de training data}{Deck oficial de Loubna Ben Allal, página 45}

\begin{warningbox}{La tasa de filtrado no es un quality metric}
Eliminar el 90\% puede limpiar duplicates, pero también puede borrar rare languages; conservar el 90\% puede mantener la diversity, pero también puede conservar spam. Hay que reportar los tipos eliminados, la source/language retention, la token distribution y el downstream effect.
\end{warningbox}

\subsection{Language selection y per-language thresholds}

v1 selecciona unos 86 de 358 languages y excluye configs y lenguajes sin mantenimiento; v2 se amplía a 600+. Estadísticas como la longitud media de línea, el tamaño de archivo o el comment ratio no pueden compartir un mismo umbral, porque cada lenguaje tiene idioms y generated-code patterns distintos.

\lecturefigure{slide-46.jpg}{The Stack v1: selección de lenguajes, muestreo comunitario y per-language filtering}{Deck oficial de Loubna Ben Allal, página 46}

\begin{knowledgebox}{La forma concreta del humano en el ciclo}
La comunidad BigCode revisa unas 100 muestras por extension, registra el código real, los generated/vendor data, las longitudes y los formatos anómalos, y a partir de ello deriva heuristics. No se trata de que los voluntarios aprueben archivo por archivo, sino de calibrar scalable rules con domain knowledge.
\end{knowledgebox}

\teachervoice{00:30:00--00:33:20, el ponente enfatiza que un mismo average-line-length threshold no puede usarse para todos los lenguajes; la inspection de la comunidad proporcionó el fundamento para diseñar umbrales por lenguaje.}

\subsection{Near-dedup: por qué MinHash y LSH se adaptan a TB-scale code}

El exact dedup solo elimina archivos idénticos; el near-dedup debe identificar cambios de nombre, cambios de formato o modificaciones locales. Si cada archivo se representa como un conjunto de shingles $A,B$, la Jaccard similarity es
\[
J(A,B)=\frac{|A\cap B|}{|A\cup B|}.
\]
MinHash cumple que la probabilidad de colisión del hash es aproximadamente $\Pr[h(A)=h(B)]=J(A,B)$; después, LSH (Locality-Sensitive Hashing) coloca las signatures similares en cubetas candidatas, lo que evita la comparación de todos contra todos.

\lecturefigure{slide-47.jpg}{Near-deduplication: MinHash+LSH y la ganancia en la ablation de StarCoder}{Deck oficial de Loubna Ben Allal, página 47}

\begin{importantbox}{Por qué el near-dedup es un filter de alto valor}
Es independiente del lenguaje, fácil de paralelizar y reduce directamente el entrenamiento repetido y la memorization. La ablation de la slide muestra que pass@1 mejora claramente en varios code subsets, por lo que el equipo adoptó un strong dedup; esto es más confiable que la heuristic de stars/comments sin validar.
\end{importantbox}

\teachervoice{00:33:20--00:35:20, el ponente dice que el near-dedup es el filter con mayor efecto y que no requiere modificarse por separado para cada uno de los 86 lenguajes; cumple a la vez con la ganancia y la escalabilidad.}

\subsection{PII, secrets y benchmark decontamination}

Ni siquiera el secret scanning integrado de GitHub garantiza que el corpus esté libre de emails, names, passwords y keys. BigCode primero anotó manualmente PII spans y luego entrenó el StarPII NER model para escanear la totalidad de los datos; en clase se informó que este paso consumió unas 800 GPU hours. Después hay que retirar también los evaluation benchmarks para evitar un inflated score.

\lecturefigure{slide-48.jpg}{The Stack: PII removal y benchmark decontamination}{Deck oficial de Loubna Ben Allal, página 48}

\begin{knowledgebox}{Primera definición: PII y decontamination}
La PII (Personally Identifiable Information) incluye names, emails, phones y otros datos que permiten identificar a una persona; los secrets incluyen además tokens/passwords/keys. La \textbf{Decontamination} consiste en retirar de los datos de entrenamiento los benchmark items y sus versiones aproximadas o derivadas, para proteger la evaluation independence.
\end{knowledgebox}

\begin{warningbox}{Un Detector solo reduce el riesgo; no puede certificar «cero PII»}
El NER tiene false negatives, las cadenas de código pueden estar ofuscadas y un commit histórico puede filtrar contenido eliminado. Al publicar se necesitan además sample audits, un reporting channel, gating e incident response.
\end{warningbox}

\teachervoice{00:34:00--00:36:20, el ponente describe con mucho detalle la annotation, el NER y el costo de 800 GPU-hour del tratamiento de la PII, y también recuerda que el evaluation set debe retirarse del training corpus.}

\subsection{Formatting: convertir el repository context en señal de entrenamiento}

Después del filtrado aún hay que decidir el sequence format. StarCoder agrega metadata tokens de repository/name/file antes del code; los issues incluyen title, comment y user roles. El formato permite que el modelo aprenda los límites del contexto y también define una interfaz que puede controlarse durante la inferencia.

\lecturefigure{slide-49.jpg}{StarCoder data formatting: metadata de repository, file, stars e issue}{Deck oficial de Loubna Ben Allal, página 49}

\lecturefigure{slide-50.jpg}{Formato estructurado de Git commits y Jupyter notebooks}{Deck oficial de Loubna Ben Allal, página 50}

Para el fill-in-the-middle (FIM), el documento se divide en prefix $P$, middle $M$ y suffix $S$, y la secuencia de entrenamiento puede reordenarse como
\[
\langle\text{PRE}\rangle P\langle\text{SUF}\rangle S
\langle\text{MID}\rangle M,
\]
se sigue usando el next-token loss, pero el modelo aprende a generar el code intermedio cuando conoce el contexto anterior y el posterior.

\begin{warningbox}{Los Metadata tokens también pueden convertirse en un shortcut}
Si las stars, el repository name o la user identity están muy correlacionados con las etiquetas de calidad, el modelo puede memorizar el source bias. Hay que evaluar la capacidad tras retirar los metadata, el privacy risk y la generalization entre repositorios.
\end{warningbox}

\subsection{StarCoder2 mixture: cada source aporta tareas distintas}

La tabla de datos de entrenamiento enumera The Stack, GitHub issues, pull requests, notebooks, documentation, math/code datasets, etc., e indica si se aplica FIM, las epochs y el sample weight. Al leer la tabla no basta con mirar el porcentaje de tokens; hay que preguntarse qué behavior enseña cada tipo de source.

\lecturefigure{slide-51.jpg}{StarCoder2 training-data mixture y source-specific settings}{Deck oficial de Loubna Ben Allal, página 51}

Si la source $i$ tiene $T_i$ tokens y un sampling weight $w_i$, la probabilidad real de muestreo es
\[
p_i=\frac{w_iT_i}{\sum_j w_jT_j}.
\]
Una source pequeña puede someterse a oversample con $w_i>1$, pero la exposure repetida también aumenta la memorization, por lo que debe calibrarse con validation/evaluation.

\begin{knowledgebox}{Source-to-behavior mapping}
Los Code files enseñan syntax/algorithms; los issues enseñan problem description y discussion; los commits/PRs enseñan edit rationale; los notebooks enseñan el entrelazado code--output--markdown; la documentation enseña API usage; los math datasets enseñan symbolic/reasoning. La Mixture es diseño de capacidades, no solo concatenación.
\end{knowledgebox}

\subsection{Tooling: un pipeline de datos reproducible no puede esconderse en un notebook de un solo uso}

La última página enumera datasets, datatrove, datatrove-like distributed pipelines, tokenizers/training code y el BigCode processing repo. El TB-scale filtering requiere streaming, sharding (dividir archivos o muestras en fragmentos que distintos workers pueden procesar de forma independiente), checkpointing, deterministic transforms y dataset versioning.

\lecturefigure{slide-52.jpg}{Tooling de filtrado de datos y entrenamiento: datasets, datatrove, frameworks de entrenamiento y BigCode code}{Deck oficial de Loubna Ben Allal, página 52}

\begin{lstlisting}[caption={Contrato por etapas de un pipeline de datos reproducible}]
discover_sources_and_record_provenance()
normalize_and_language_classify()
apply_quality_filters_with_versioned_thresholds()
exact_and_near_deduplicate()
detect_pii_secrets_and_benchmark_overlap()
format_sequences_and_sample_mixture()
publish_manifest_hashes_stats_and_opt_out_state()
\end{lstlisting}

\subsection{Resumen de la sección}

La data advantage de StarCoder no proviene de un filter mágico, sino de la combinación de source selection, per-language inspection, near-dedup, PII/decontamination, structured formatting, mixture y reproducible tooling. Cada paso tiene costos, modos de falla y evidence.
\section{El ecosistema de modelos de código: de Copilot a la colaboración BigCode}

Una vez completo el pipeline de datos, la clase vuelve a la historia y al ecosistema de los code LLM. El cambio central es que Codex/Copilot demostró que la next-token prediction basta para constituir un baseline fuerte; después, la comunidad combinó datos, modelos y fine-tuning públicos en una gran cantidad de code assistants.

\subsection{2021: el código puede tratarse como texto con next-token prediction}

El code modeling temprano solía usar de forma explícita AST, program analysis o feature engineering; Codex mostró que, si se alimenta un Transformer grande con grandes cantidades de code, también puede aprender completion. Los métodos estructurados siguen teniendo valor, pero ya no son un requisito previo para llegar a una code generation fuerte.

\lecturefigure{slide-55.jpg}{GitHub Copilot 2021: el punto de partida como producto del next-token code modeling}{Deck oficial de Loubna Ben Allal, página 55}

\teachervoice{00:38:30--00:40:00, la ponente resume el significado histórico de Codex como «code can be treated like text»: no es que el código carezca de estructura, sino que un sequence baseline fuerte reduce mucho la barrera que imponían las características diseñadas a mano.}

\subsection{2024: más de 1.7K open code models y el panorama de los leaderboards}

En el Hub aparecen muchos modelos code-pretrained o con una code mixture; al mismo tiempo, los base models y los instruction-tuned models mejoran con rapidez en tablas de clasificación del tipo de HumanEval. La cantidad indica que el ecosystem está activo, pero no equivale a 1,700 pretraining runs independientes de alta calidad: muchos son fine-tunes, quantizations o derivatives.

\lecturefigure{slide-56.jpg}{Panorama de 2024: más de 1.7K open models trained on code}{Deck oficial de Loubna Ben Allal, página 56}

\lecturefigure{slide-57.jpg}{Panorama de las tablas de clasificación de code base models e instruction-tuned models fuertes}{Deck oficial de Loubna Ben Allal, página 57}

\begin{warningbox}{80\% en HumanEval no es «80\% de capacidad de programación»}
HumanEval consta de unas 164 Python function-completion tasks, evaluadas con unit tests. La puntuación depende del prompt, el sampling, pass@k, la contamination y el instruction tuning; no puede extrapolarse a repository editing, debugging, seguridad ni ingeniería de software con múltiples archivos.
\end{warningbox}

\subsection{Open code LLM landscape: provider, size y objective}

Meta Code Llama, BigCode StarCoder, DeepSeek Coder y otros forman distintas families; además están CodeGen, StableCode, CodeQwen y Granite. Al compararlos, conviene distinguir base/instruct, context, FIM, license, languages, training date y benchmark protocol.

\lecturefigure{slide-58.jpg}{Open code LLM landscape: Meta, BigCode, DeepSeek y otras families}{Deck oficial de Loubna Ben Allal, página 58}

\begin{knowledgebox}{Términos clave de los code models (primera aparición)}
Un \textbf{base model} aprende principalmente next-token/FIM; un \textbf{instruction-tuned model} aprende a responder instrucciones; \textbf{MQA} hace que todas las query heads compartan un solo conjunto de KV; \textbf{GQA} hace que varias query heads compartan un conjunto de KV, lo que por lo general es un punto intermedio entre calidad y KV-cache.
\end{knowledgebox}

\teachervoice{00:55:00--00:57:20, en la sesión de Q\&A la ponente dice que el entrenamiento con code y con text es similar en general, pero que un code assistant da más peso a long context, a la inference efficiency de MQA/GQA y a smaller models adecuados para el IDE (IDE-friendly).}

\subsection{BigCode: el objetivo de una colaboración de ciencia abierta}

BigCode no solo publica weights: hace que más de 1,000 investigadores e ingenieros participen en Slack y hace públicos el data processing, el training code, los models y una friendly license. El valor está en poder revisar el decision trail, no en tomar el collaboration size como garantía de calidad.

\lecturefigure{slide-59.jpg}{BigCode open-scientific collaboration: data, processing, models y comunidad transparentes}{Deck oficial de Loubna Ben Allal, página 59}

\lecturefigure{slide-60.jpg}{BigCode ecosystem: The Stack, StarCoder y los fine-tunes de la comunidad}{Deck oficial de Loubna Ben Allal, página 60}

\begin{importantbox}{La cadena de reutilización del ecosystem}
The Stack lo usan varios code models; a su vez, los checkpoints de StarCoder se vuelven la base de fine-tunes como StarChat y WizardCoder. La reutilización amplía el impacto, pero también significa que las decisiones de upstream sobre data quality, license y contamination se propagan hacia downstream.
\end{importantbox}

\teachervoice{00:40:00--00:43:10, la ponente dice que la razón directa para iniciar BigCode fue la full data transparency: la comunidad puede inspeccionar los datos (inspect data), reutilizar el processing code y revisar los modelos, en lugar de solo descargar weights.}

\subsection{Resumen de la sección}

El ecosistema de los code LLM pasó con rapidez del Copilot de código cerrado a datasets, base models y fine-tunes públicos. La siguiente pregunta ya no es «si hay modelos», sino si un release es verdaderamente abierto, transparente, reproducible y responsable.
\section{Open and Responsible: los pesos abiertos son solo una cuarta parte}

Esta sección coloca la data governance en el centro del release de un modelo. La apertura comprende, como mínimo, datos accesibles, modelos accesibles, investigación reproducible y documentación completa; la responsibility exige además opt-out, PII, license, evaluation y limitation reporting.

\subsection{Cuatro niveles de apertura: Data, Models, Research, Documentation}

La slide divide la apertura en cuatro columnas. Los datos requieren inspection/opt-out/PII; los modelos requieren weights, training code y evaluation interface; la investigación requiere processing/training/evaluation reproducibility; la documentación requiere dataset/model cards, governance, architecture, cost y limitations.

\lecturefigure{slide-61.jpg}{La checklist de cuatro niveles de Open and Responsible Research on LLMs}{Deck oficial de Loubna Ben Allal, página 61}

\begin{importantbox}{Definición mínima de un release abierto}
\begin{enumerate}
  \item Poder saber qué data y qué transformation utilizó el modelo;
  \item Poder obtener, o al menos auditar, weights/code/config;
  \item Poder reproducir el training/evaluation path principal;
  \item Poder consultar license, limitations, PII/opt-out y cost.
\end{enumerate}
Contar solo con downloadable weights no basta para responder estas preguntas.
\end{importantbox}

\teachervoice{00:43:10--00:45:10, el ponente afirma explícitamente que un open release no consiste en «release model weights and stop there», sino en la combinación de data tools, opt-out, PII, reproducibility, evaluation tools y technical reports.}

\subsection{SantaCoder, StarCoder, StarCoder2: cómo itera el proyecto}

SantaCoder sirvió como ablation platform de alrededor de 1.1B; StarCoder escaló a 15B/alrededor de 1T tokens; StarCoder2 ofrece versiones de 3B/7B/15B y utiliza un corpus más grande y con más lenguajes. La slide muestra a la vez la disminución de la proporción de Python source y la continuidad de la transparencia y el acceso abierto.

\lecturefigure{slide-62.jpg}{De SantaCoder a StarCoder2: escala, lenguajes, proporción de Python y apertura}{Deck oficial de Loubna Ben Allal, página 62}

\begin{knowledgebox}{Por qué empezar con SantaCoder}
Una TB-scale data rule no puede probarse repetidamente de forma directa en un full run de 15B. El modelo pequeño ofrece una ablation platform para filter/format/objective; sin embargo, la extrapolation aún debe validarse en modelos más grandes, y no puede suponerse que todo data effect sea completamente scale-invariant.
\end{knowledgebox}

\subsection{Transparency Index: un indicador externo mide la divulgación, no toda la capacidad}

El Stanford Foundation Model Transparency Index asigna puntuaciones según dimensiones de divulgación como data, labor, compute, risks y downstream. En la instantánea presentada en clase, StarCoder aparece en los primeros lugares, lo que indica que el esfuerzo de documentación/transparencia fue captado por un marco externo; esto no demuestra que el modelo sea el más preciso ni el más seguro, ni que las controversias legales estén resueltas.

\lecturefigure{slide-63.jpg}{Resultados del Foundation Model Transparency Index al momento de la clase}{Deck oficial de Loubna Ben Allal, página 63}

\begin{warningbox}{Transparency y responsibility están relacionadas, pero no son equivalentes}
Divulgar malos resultados es más transparente que ocultarlos, pero aun así puede ser irresponsable; sin una divulgación completa puede haber controles internos, pero el público no puede verificarlos. El Index mide la información visible y no sustituye una independent audit ni un legal judgment.
\end{warningbox}

\subsection{Resumen de la sección}

Un responsible open model requiere publicar en conjunto data, weights, code, evaluation, documentation y governance. El valor de BigCode reside en convertir estos niveles en objetivos del proyecto; los indicadores de transparencia aportan una revisión externa, pero no sustituyen la evaluación de capacidades y riesgos.
\section{Evaluation y Tooling: primero prevenir la Contamination, después hablar de Leaderboards}

Que termine el preentrenamiento no significa que los resultados sean confiables. Los benchmarks de código tienden especialmente a aparecer en GitHub, en tutoriales y en instruction data; además, los leaderboards públicos incentivan la optimización repetida. Por ello, el curso propone usar múltiples benchmarks, decontamination, membership test y time-split evaluation.

\subsection{Múltiples benchmarks: una sola puntuación no describe un code model}

La primera tabla compara al mismo tiempo code completion, code reasoning, data science, math y repository tasks. Un modelo de mayor size suele ser más fuerte, pero el orden de los modelos puede variar entre tasks. La evaluación debe ir ligada a model version, prompt, sampling, tests y execution environment.

\lecturefigure{slide-64.jpg}{Comparación en múltiples benchmarks de StarCoder/CodeLlama/DeepSeek y otros modelos}{deck oficial de Loubna Ben Allal, página 64}

\lecturefigure{slide-65.jpg}{Resultados desglosados por tarea de StarCoder2 3B/7B/15B}{deck oficial de Loubna Ben Allal, página 65}

Si para cada problema se muestrean $n$ solutions y $c$ de ellas pasan los tests, el estimador insesgado habitual de pass@$k$ es
\[
\operatorname{pass@}k=1-\frac{\binom{n-c}{k}}{\binom{n}{k}}.
\]
pass@1 se aproxima más a una sola completion; un $k$ mayor mide que al menos uno de varios intentos tenga éxito, por lo que los valores no pueden compararse directamente.

\begin{warningbox}{Los unit tests solo cubren el comportamiento que se escribió}
Pasar los tests no garantiza efficiency, security, style, ausencia de undefined behavior, dependency correctness ni el manejo de hidden edge cases. Si los tests son demasiado débiles, el modelo puede generar código que parece correcto pero no es confiable.
\end{warningbox}

\teachervoice{00:44:20--00:46:20, el ponente recomienda usar en el release tantos benchmarks como sea posible: alguno puede tener contamination, y las múltiples tareas también revelan la distribución del comportamiento del modelo, en lugar de perseguir solo un headline individual.}

\subsection{Autocomplete y membership test: el lado del despliegue también debe hacer attribution}

BigCode también publicó una VS Code extension y un membership test. Este último intenta determinar si un generated snippet se parece a los training data y advierte al usuario sobre su posible origen. Es un provenance aid, no una memorization proof concluyente.

\lecturefigure{slide-66.jpg}{Code tooling: autocomplete extension y training-data membership test}{deck oficial de Loubna Ben Allal, página 66}

\begin{knowledgebox}{Límites del membership test}
El exact match tiene alta precisión pero pasa por alto las reescrituras; el fuzzy match marca common idioms como falsos positivos; un training corpus incompleto o de otra versión produce false negatives. La salida debe ser una similitud y unas fuentes candidatas, no redactarse como una conclusión legal.
\end{knowledgebox}

\teachervoice{00:45:00--00:47:00, el ponente ubica el membership test dentro de los esfuerzos de code attribution: cuando el código generado puede provenir de ejemplos de entrenamiento, se debe dar al autor o al usuario una advertencia que pueda verificar.}

\subsection{Customize code model: cómo llevar los principios del preentrenamiento a una sola GPU}

El Personal Copilot workflow sigue siendo collect--filter--dedup--format--train--evaluate; lo único que cambia es que el full pretraining se sustituye por repository-specific continued pretraining o fine-tuning. LoRA expresa la actualización de los pesos como una matriz de bajo rango
\[
W'=W+BA,\qquad A\in\mathbb{R}^{r\times d},\ B\in\mathbb{R}^{d'\times r},\ r\ll d,
\]
y solo entrena $A,B$, lo que reduce notablemente la memory de los optimizer states y de los gradientes.

\lecturefigure{slide-67.jpg}{Customize Code Models: entrenar un personal Copilot sobre un codebase privado}{deck oficial de Loubna Ben Allal, página 67}

\begin{lstlisting}[caption={Ruta ejecutable para un code assistant con presupuesto reducido}]
select_small_strong_base_with_code_and_long_context()
collect_only_authorized_repository_data()
deduplicate_and_split_by_repository_or_time()
format_fim_and_repository_context_examples()
fine_tune_with_peft_or_lora_on_one_gpu()
evaluate_execution_latency_leakage_and_regressions()
\end{lstlisting}

\teachervoice{00:57:20--00:59:10, en la sesión de Q\&A el ponente da recomendaciones para una sola GPU: elegir una base fuerte, curated data, un runtime quantized/on-device y PEFT, en lugar de intentar preentrenar desde cero.}

\subsection{Leaderboards: la interfaz es cómoda, pero el protocol importa más}

Los leaderboards públicos agregan model scores y ayudan a identificar candidatos; pero detrás de cada row hay model date, prompt, temperature, test harness, contamination policy y submission rules. La captura de pantalla es un punto de entrada, no un registro de experimento reproducible.

\lecturefigure{slide-68.jpg}{Code leaderboards: punto de entrada público para comparar modelos}{deck oficial de Loubna Ben Allal, página 68}

\begin{knowledgebox}{Orden de uso de un leaderboard}
Primero revisa la benchmark version y el cutoff; después, la release/training date del modelo; confirma base/instruct, pass@1/pass@k, execution sandbox y prompt; solo al final compara las puntuaciones. Si el protocol es distinto, la clasificación carece de interpretabilidad.
\end{knowledgebox}

\subsection{HumanEval overfitting: el instruction tuning también contamina}

Aunque el pretraining elimine HumanEval, los instruction data pueden contener ejercicios parecidos del tipo «implementar una función y pasar los tests». El modelo puede aprender el benchmark style o haber visto problemas semánticamente equivalentes, lo que infla el public score respecto de problemas realmente nuevos.

\lecturefigure{slide-69.jpg}{Potential overfitting in HumanEval: brecha entre el benchmark público y la generalización real}{deck oficial de Loubna Ben Allal, página 69}

Si la similarity entre el ejemplo de entrenamiento $x$ y el benchmark $b$ cumple $s(x,b)>\tau$, puede marcarse un posible overlap; sin embargo, un lexical detector no necesariamente capta la semantic equivalence. La decontamination debe combinar exact/fuzzy/code-structure/time metadata y reportar los umbrales.

\begin{warningbox}{Los instruction data son la segunda vía de entrada de la contamination}
No basta con limpiar el pretraining corpus. Los SFT, los RLHF prompts, los synthetic exercises, las benchmark explanations y las GitHub solutions pueden filtrar el test pattern; cada stage requiere un provenance scan.
\end{warningbox}

\subsection{LiveCodeBench: aislamiento temporal mediante la release date}

LiveCodeBench recopila periódicamente nuevos contest problems y evalúa solo con los problemas posteriores al release del modelo, lo que reduce la probabilidad de que este los haya visto durante el entrenamiento. También compara el full-history score con el post-cutoff score; si la diferencia es grande, puede indicar overfit a benchmarks antiguos.

\lecturefigure{slide-70.jpg}{LiveCodeBench: contamination-resistant evaluation con post-release problems}{deck oficial de Loubna Ben Allal, página 70}

La time-split puede definirse como
\[
\mathcal{B}_{\text{valid}}(m)=\{b:\ t_{\text{problem}}(b)>t_{\text{release}}(m)\}.
\]
Reduce la fuga temporal directa, pero el modelo puede basarse en un checkpoint no publicado, en problemas que aparecieron antes en la web o en templates similares, por lo que se requiere cautela.

\teachervoice{00:49:00--00:52:40, el ponente presenta LiveCodeBench como una solución importante: usar problemas que aparecieron después de la publicación del modelo y observar si algunos modelos mantienen su desempeño en el uncontaminated slice.}

\subsection{Resumen de la sección}

Una code evaluation confiable requiere múltiples tareas, tests ejecutables, un pass@k explícito, decontamination de datos y temporal, herramientas de membership/attribution y un repeated protocol. El leaderboard sirve para orientarse; la time-split y el source audit son la defensa real contra la contamination.
\section{Q\&A: Domain, Tokenizer, Fine-tuning y Dataset Release}

La sesión de Q\&A lleva la lecture recipe a escenarios límite: si los multimodal data pueden sustituir al text, si el code requiere una architecture especial, cómo trabajar con una sola GPU, si el scaling se mantiene entre domains, qué diferencias hay entre tokenizers, si el fine-tuning usa el mismo filter y qué gobernanza adicional requiere publicar un conjunto de datos grande.

\subsection{Code y text: training similar, prioridades distintas}

La ponente considera que la architecture/training es similar en general; un code model puede seguir usando un decoder Transformer. Las diferencias están más bien en el long repository context, FIM, la fast IDE inference, MQA/GQA y el data formatting. El code tokenizer sigue usando BPE, pero conserva el number splitting y, sobre la code mixture, se revisa si hay tokens under/overrepresented.

\begin{knowledgebox}{Primera definición de tokenizer y sanity checks}
BPE (Byte Pair Encoding) fusiona de forma repetida los token pairs más frecuentes. Un tokenizer de código debe revisar whitespace/indentation, operators, numbers, Unicode, paths y language coverage. Un token-per-character promedio demasiado alto aumenta el context cost, mientras que unos whole-identifier tokens demasiado grandes perjudican la generalización composicional.
\end{knowledgebox}

\teachervoice{00:58:20--01:00:00: la ponente dice que el code tokenizer es muy parecido al text tokenizer; lo importante es el number splitting, entrenarlo sobre la mixture real y revisar los outlier tokens. El código también contiene mucho Markdown/English.}

\subsection{El preentrenamiento y el fine-tuning usan filters de distinta intensidad}

El preentrenamiento necesita breadth: un star filter demasiado estricto elimina demasiados datos. El fine-tuning se orienta a un repository/language específico, por lo que puede seleccionar muestras de alta calidad con mayor rigor. El resultado de LIMA con 1,000 instructions indica que en la etapa de alignment la quality puede importar más que la quantity, pero eso no significa que 1,000 muestras basten para cualquier tarea.

\begin{warningbox}{No extrapolar LIMA a «pocos datos siempre bastan»}
LIMA usa una pretrained base sólida y un alignment setting específico. Si el domain knowledge no está en la base, el fine-tuning sigue necesitando una cobertura suficiente; con pocos datos también es más fácil sobreajustar el estilo, filtrar datos de evaluación o carecer de rare cases.
\end{warningbox}

\teachervoice{00:59:20--01:00:20: la ponente contrasta ambos casos: en el preentrenamiento no se puede recortar a la ligera hasta perder breadth; el fine-tuning necesita menos datos, por lo que se puede dedicar más esfuerzo a un filtrado estricto.}

\subsection{Multimodal, domain scaling y respuestas inciertas}

Cuando se le pregunta si image/video reducen la necesidad de text-only data, la ponente afirma con claridad que no lo ha probado y que no puede dar una proporción; solo observó que los VLM de ese momento todavía usaban mucho text. Ante la pregunta sobre domains como el medical, considera que el scaling optimum podría cambiar y que la evidence disponible es insuficiente.

\begin{importantbox}{Conservar el «no sé» también es teacher voice}
Una nota de alta calidad no debe completar por la ponente una respuesta que ella no dio. La conclusión aplicable aquí es: para una nueva modality/domain hay que repetir la scaling/data-mixture ablation, en lugar de citar la generic Chinchilla ratio.
\end{importantbox}

\teachervoice{00:52:40--00:55:00: sobre la multimodal data substitution, la ponente responde directamente «no lo hemos probado, no puedo responderlo realmente» y solo ofrece un juicio basado en observaciones.}

\subsection{Publicar conjuntos de datos grandes: diseñar juntos la parte técnica, la gobernanza y el control de acceso}

Antes de publicar se necesitan license/copyright review, opt-out, PII/secrets detection, dataset card, processing code, statistics, gating de los sensitive subsets y una versioned removal list. Hub/gated access mejora la accesibilidad y el control de riesgos, pero no constituye una exención legal.

\begin{lstlisting}[caption={Checklist previa a la publicación de un conjunto de datos de entrenamiento grande}]
record_source_license_timestamp_and_transform_lineage()
publish_filter_dedup_pii_and_decontamination_versions()
provide_inspection_opt_out_and_removal_workflows()
gate_sensitive_annotations_or_high_risk_subsets()
release_statistics_cards_tools_and_known_limitations()
version_updates_and_propagate_removals_to_future_models()
\end{lstlisting}

\teachervoice{01:00:00--01:01:30: la ponente recomienda considerar al mismo tiempo license, copyright, opt-out, herramientas y documentación; para sensitive subsets como los PII annotation data, se puede usar gated access.}

\subsection{Resumen de la sección}

Un Code LLM no es una especie completamente distinta: reutiliza la recipe de los modelos de lenguaje, pero con prioridades diferentes en context, inference, format y evaluation. Con un presupuesto limitado conviene hacer curated fine-tuning; los nuevos domains/modalities exigen volver a medir el scaling, y la publicación de datos debe tratar los flujos de trabajo de gobernanza como una funcionalidad del producto.
\section{Síntesis y ampliación}

Lo que esta clase muestra en realidad es un data-to-model control loop: primero se usan el scaling y el lifecycle cost para decidir la escala objetivo; luego se eligen las fuentes y la provenance; después se buscan los filtros con ablation, se atienden opt-out/PII/license con governance, se define la tarea con formatting/mixture y, por último, se revisa el modelo con contamination-aware evaluation. El valor de StarCoder no está solo en el code score: está en que toda esta cadena puede auditarse.

\subsection{Cadena causal de extremo a extremo}

\begin{importantbox}{De la data decision al model behavior}
\begin{enumerate}
  \item \textbf{Source} determina el conocimiento, las licencias y la cobertura de la cola larga;
  \item \textbf{Filter/dedup} determina la calidad, la diversity y la memorization;
  \item \textbf{Formatting/mixture} determina las tareas y el contexto que ve el modelo;
  \item \textbf{Training scale} determina la capacidad y la exposure;
  \item \textbf{Evaluation} determina si podemos distinguir la generalización de la contamination.
\end{enumerate}
Si cualquiera de estas capas no es visible, el benchmark final es difícil de interpretar.
\end{importantbox}

\subsection{Los cinco errores más comunes}

\begin{warningbox}{Del headline a una conclusión falsa}
No presentes la Chinchilla ratio como una constante válida entre domains; no uses stars/comments como quality label; no presentes un PII detector como prueba de riesgo cero; no presentes la puntuación de HumanEval como capacidad de ingeniería de software; no presentes los open weights como openness completa. En el curso, los resultados negativos y las limitaciones de gobernanza importan tanto como los resultados positivos.
\end{warningbox}

\subsection{Lista de verificación para reproducir los experimentos}

\begin{lstlisting}[caption={Lista unificada de reproducción para proyectos de preentrenamiento}]
pin_model_tokenizer_optimizer_and_code_commit()
hash_raw_filtered_and_formatted_dataset_versions()
report_source_mixture_repeats_and_effective_tokens()
publish_filter_ablations_seeds_and_negative_results()
document_opt_out_pii_license_and_decontamination()
freeze_evaluation_prompts_harnesses_cutoffs_and_dates()
report_train_and_inference_compute_memory_latency_cost()
\end{lstlisting}

\subsection{Preguntas de autoevaluación}

\begin{multicols}{2}
\footnotesize
\begin{enumerate}
  \setlength{\itemsep}{0pt}
  \item ¿Por qué los open weights no equivalen a un transparent release completo?
  \item ¿En qué difieren Kaplan y Chinchilla respecto de la model/data allocation?
  \item ¿Por qué lo compute-optimal no es necesariamente lifecycle-optimal?
  \item ¿Qué diferencia hay entre raw bytes, unique tokens y training tokens?
  \item ¿Por qué no se puede tokenizar Common Crawl directamente y entrenar con él?
  \item ¿Qué data lineage requiere el opt-out de The Stack v1?
  \item ¿Qué cambia en v2 el uso de Software Heritage?
  \item En los synthetic data, ¿qué controla cada uno: el seed, el generator y el verifier?
  \item ¿Cómo aumenta Cosmopedia la diversidad de lo generado?
  \item ¿Por qué el repository stars filter produjo la peor ablation?
  \item ¿Para qué sirven varios seeds en un filter experiment?
  \item ¿Por qué MinHash/LSH es adecuado para la near-dedup?
  \item ¿Por qué una strong dedup puede mejorar un code benchmark?
  \item ¿Qué diferencia hay entre PII detection y secret scanning?
  \item ¿Por qué la decontamination debe abarcar los SFT data?
  \item ¿Cómo cambian los metadatos de repository/file la tarea de entrenamiento?
  \item ¿Cómo se implementa FIM con un next-token objective?
  \item ¿Cómo cambia el source sampling weight la effective mixture?
  \item ¿Qué optimizan principalmente MQA y GQA?
  \item ¿Por qué la openness de BigCode va más allá de publicar los pesos?
  \item ¿Pueden compararse directamente pass@1 y pass@k?
  \item ¿Por qué un membership test no puede demostrar memorization?
  \item ¿Qué tipo de fuga reduce la time split de LiveCodeBench?
  \item Para construir un code assistant con una sola GPU, ¿en qué conviene invertir primero?
\end{enumerate}
\end{multicols}

\subsection{Lecturas adicionales}

\begin{multicols}{2}
\footnotesize
\begin{itemize}
  \setlength{\itemsep}{0.15em}
  \item Hoffmann et al., \href{https://arxiv.org/abs/2203.15556}{Training Compute-Optimal Large Language Models}: Chinchilla scaling.
  \item Muennighoff et al., \href{https://arxiv.org/abs/2305.16264}{Scaling Data-Constrained Language Models}: entrenamiento con repetición cuando los unique data son limitados.
  \item Kocetkov et al., \href{https://arxiv.org/abs/2211.15533}{The Stack}; Lozhkov et al., \href{https://arxiv.org/abs/2402.19173}{StarCoder 2 and The Stack v2}.
  \item Penedo et al., \href{https://arxiv.org/abs/2306.01116}{RefinedWeb} y \href{https://arxiv.org/abs/2406.17557}{FineWeb}: web filtering y ablations.
  \item Lee et al., \href{https://arxiv.org/abs/2107.06499}{Deduplicating Training Data Makes Language Models Better}.
  \item Li et al., \href{https://arxiv.org/abs/2305.06161}{StarCoder: may the source be with you!}: modelo, FIM, gobernanza y evaluación.
  \item Brown et al., \href{https://arxiv.org/abs/2107.03374}{HumanEval}; Jain et al., \href{https://arxiv.org/abs/2403.07974}{LiveCodeBench}: evaluación por execution y con separación temporal.
  \item Zhou et al., \href{https://arxiv.org/abs/2305.11206}{LIMA}; Hu et al., \href{https://arxiv.org/abs/2106.09685}{LoRA}: alignment con pocos datos de alta calidad y parameter-efficient fine-tuning.
\end{itemize}
\end{multicols}

\end{document}
